\documentclass[11pt]{article}


\setlength{\parindent}{0pt}
\usepackage{xltxtra}
\usepackage{hyperref}
\hypersetup{hidelinks}
\usepackage{url}
\urlstyle{tt}
\usepackage{xcolor}
\definecolor{CVBlue}{RGB}{23,110,191}
\usepackage{calc}
\usepackage{graphicx}
\usepackage{tikz}
\usetikzlibrary{calc}
\usepackage{fontspec}
\usepackage{xeCJK}
\usepackage{enumitem}
\CJKsetecglue{} %% 取消中文与数字之间的间隙


%% 主文档字体设置
\setmainfont[
    Path = fonts/Main/,
    Extension = .otf,
    BoldFont = texgyretermes-bold.otf, % 加粗字体
]{texgyretermes-regular.otf} % 正文字体

% 中文字体设置
\setCJKmainfont[
    Path = fonts/hansans/,
    Extension = .ttf,
    BoldFont = NotoSansSC-Bold.ttf, % 加粗字体
]{NotoSansSC-Regular.ttf} % 正文字体


\usepackage{relsize}
\usepackage{xspace}

% 使用 fontawesome（部分图标）
\usepackage{fontawesome}

% A4纸，上下左右边距
\usepackage[
    a4paper,
    left=1.2cm,
    right=1.2cm,
    top=1.5cm,
    bottom=1cm,
    nohead
]{geometry}

\renewcommand{\baselinestretch}{1.5} % 行间距设为1.5

\usepackage{titlesec}
\usepackage{enumitem}
\setlist{noitemsep} % 取消列表项间的额外间距
%\setlist{nosep} % 取消所有垂直间距
\setlist[itemize]{topsep=0.25em, leftmargin=*}
\setlist[enumerate]{topsep=0.25em, leftmargin=*}

% --- 用于控制【不同项目之间】的垂直距离 ---
\newlength{\interProjectSpacing}
\setlength{\interProjectSpacing}{0.9em} % <--- 在此调整项目之间的距离
\newcommand{\projectsep}{\vspace{\interProjectSpacing}}

% --- 用于控制【项目标题】与下方【项目描述】的距离 ---
\newlength{\intraProjectTitleSep}
\setlength{\intraProjectTitleSep}{0.4em} % <--- 在此调整标题和描述的距离
\newcommand{\titlebreak}{\\[\intraProjectTitleSep]}

% --- 用于控制【项目描述】与下方【要点列表】的距离 ---
\newlength{\intraProjectListTopSep}
\setlength{\intraProjectListTopSep}{0.2em} % <--- 在此调整描述和列表的距离

% =======================================================================


\titleformat{\section}         % 定制 \section 命令
{\large\bfseries\raggedright} % 将 section 标题设置为大号、粗体且左对齐
{}{0em}                      % 可用于为所有 section 添加前缀（如"章节..."）
{}                           % 可用于在标题前插入代码
[{\color{CVBlue}\titlerule}]  % 在标题后插入一条横线
\titlespacing*{\section}{0cm}{*1.6}{*1.2}



\begin{document}
\pagenumbering{gobble}

%%%% 利用tikz来定位照片
\begin{tikzpicture}[remember picture, overlay]
    \node[anchor = north east] at ($(current page.north east)+(-1.6cm,-0.6cm)$) {\includegraphics[height=3cm]{self_cropped.jpg}};
  \end{tikzpicture}%
  %%%% 学校Logo已移除
\vspace*{1em}
\centerline{\LARGE\bfseries{张宗森}}

\centerline{\normalsize{\faPhone\ 173-0098-8159 \quad \faEnvelopeO\ \href{mailto:17300988159@163.com}{17300988159@163.com} \quad \textcolor{CVBlue}{\faBirthdayCake}\ 2002.07 \quad \faMale\ 男}}

\vspace{1.2em}

\centerline{\parbox{0.92\textwidth}{\centering\normalsize{\textbf{\textcolor{CVBlue}{\faBullseye}\ 求职意向：}影像、多媒体、ISP、色彩、图像处理}}}

\section{\makebox[\widthof{\faGraduationCap}][c]{\color{CVBlue}\faGraduationCap}\ 教育背景}
\textbf{浙江大学} \hfill 2020.09 -- 2024.06\\
本科 \quad 光电科学与信息工程
\begin{itemize}[nosep]
    \item 相关课程：物理光学、图像处理、数字信号处理、光学设计、光电检测技术
\end{itemize}
\vspace{0.3em}
\textbf{浙江大学} \hfill 2024.09 -- 2027.06\\
硕士 \quad 光学工程 \quad （颜色科学与影像技术）
\begin{itemize}[nosep]
    \item 研究方向：HDR成像、处理与显示、Tone Mapping及其后的色彩校正、白平衡
\end{itemize}

\section{\makebox[\widthof{\faUsers}][c]{\color{CVBlue}\faUsers}\ 项目经历}

% --- 第六个项目 ---
\textbf{基于人眼色适应机制的环境光调制显示色彩补偿技术} \hfill 2025.03 -- 至今 \titlebreak
项目描述：校企合作项目。解决环境光变化带来的人眼色适应偏移问题，消除实景与拍摄回看画面存在的色彩观感差异。
\begin{itemize}[nosep, topsep=\intraProjectListTopSep]
    \item 搭建具备可调照度、色温的标准化观测实验平台，通过专业色度设备完成参数校准，依托\textbf{记忆匹配、调节法、恒常刺激法}等多类心理物理学范式采集可靠的人眼感知色貌数据。
    \item 基于 von Kries、WGM 色适应模型完成数据建模，结合人眼视觉特性优化模型结构、拓展校正模块，有效提升\textbf{色适应预测精度与泛化能力}。
    \item 研究成果可用于影像设备\textbf{自适应白点调节}功能开发，为全链路色彩统一与画质优化提供实验依据和核心算法支撑。
\end{itemize}

\projectsep

% --- 第五个项目 ---
\textbf{HDR 图像阶调映射后感知色彩补偿研究} \hfill 2025.06 -- 至今 \titlebreak
项目描述：补偿HDR图像经阶调映射后的色彩失真。该问题涉及跨亮度的主观色彩感知匹配，现有客观指标指导效果欠佳，需设计实验探索解决路径。
\begin{itemize}[nosep, topsep=\intraProjectListTopSep]
    \item 根据项目要求复现问题并进行了细致的观察分析，针对性地设计了实验方案。将问题感知表现与图像量化数据关联，拆解至\textbf{色调、饱和度}两个维度的子问题，转入感知均匀色彩空间完成色度参量对齐与主观感知匹配，分别校正了色调与饱和度偏差。
    \item 为兼顾计算效率，把优化模型转化为可快速调用的\textbf{RGB三维查找表}。经主观评估实验对比，该方法校正效果优于现有主流阶调映射色彩校正算法，项目顺利通过验收。
\end{itemize}

\projectsep

% --- 第三个项目 ---
\textbf{不可移动文物HDR图像采集与展示关键技术研究} \hfill 2024.02 -- 2024.12 \titlebreak
项目描述：山西省文物局课题。面向不可移动文物现场HDR图像采集与展示需求，建立从多曝光RAW捕获到显示设备特征化的完整技术流程。
\begin{itemize}[nosep, topsep=\intraProjectListTopSep]
    \item 从相机噪声模型出发，推导了图像信噪比随曝光参数变化的解析预测方程，并基于人眼\textbf{对比度灵敏度函数（CSF）}提出了感知导向的曝光时间选择算法，实现了HDR图像采集的效率提升。
    \item 将该曝光时间选择算法与相机控制软件 digicamcontrol 结合，编写了自动控制相机采集HDR图像的程序。
    \item 赴\textbf{云冈石窟第20窟}开展现场HDR图像采集，建立了从多曝光RAW捕获、HDR融合、阶调映射到显示设备特征化的完整采集与展示技术流程。
\end{itemize}

\projectsep

% --- 第一个项目 ---
\textbf{本科毕业设计——光学传函综合测量仪的改进设计} \hfill 2023.12 -- 2024.05 \titlebreak
项目描述：设计方案测量负焦距系统的焦距与调制传递函数（MTF）。
\begin{itemize}[nosep, topsep=\intraProjectListTopSep]
    \item 比较几种负焦距系统焦距测量方法的精度与测量范围，使用倾斜刀口法测量计算系统MTF（ESF->LSF->MTF），建模分析除目标组件外的MTF，根据级联MTF相乘关系计算得到目标组件MTF。
    \item 与合作者完成电致变焦液晶透镜（负焦距）的焦距和MTF测量，负责设计实验方案并完成测量实验，相关论文发表于\textit{ACS Applied Optical Materials}。
\end{itemize}

\section{\makebox[\widthof{\faBook}][c]{\color{CVBlue}\faBook}\ 研究成果}
\begin{itemize}[nosep]
    \item Electrically tunable liquid crystal lens with varifocal capability [J]. \textit{ACS Applied Optical Materials}, 2024. (四作)
    \item Color constancy in single-exposure imaging and multi-exposure HDR imaging [J]. \textit{Optics Express}, 2025. (四作)
\end{itemize}

\section{\makebox[\widthof{\faCogs}][c]{\color{CVBlue}\faCogs}\ 专业技能}
\begin{itemize}[nosep]
    \item \textbf{编程语言：} Python, MATLAB, C++
    \item \textbf{算法与工具：} OpenCV, MATLAB Image Processing Toolbox, Psychtoolbox, LaTeX, Linux, Git
    \item \textbf{领域专长：} ISP pipeline基础、成像与显示设备色度特征化、色貌模型、HDR成像与阶调映射、数字图像处理、颜色科学、图像质量评价、心理物理学实验设计与数据分析
    \item \textbf{深度学习：} 掌握计算机视觉主流网络架构及其在图像复原、白平衡、颜色迁移、HDR成像中的应用，熟悉 PyTorch 基础开发流程
\end{itemize}

\section{\makebox[\widthof{\faInfo}][c]{\color{CVBlue}\faInfo}\ 自我评价}
具备扎实的光学工程与颜色科学理论功底，研究方向涵盖HDR成像、色貌模型与显示技术。善于将学术研究成果转化为工程实践方案，具有较强的英文文献阅读与论文写作能力，具有多次团队合作经验，在合作中耐心负责，对待项目认真实干肯钻研，对影像品质有着源于感知实践的理解与追求，对科学技术前沿新进展保持关注，具备快速检索资源和学习新领域知识的能力。
\end{document}
